% !TeX root = ../main.tex

\section{Análisis}

\paragraph{} En primer lugar, analizaremos de forma cualitativa cómo varía la transmisión con la energía para los casos mostrados en \ref{fig:TvsE_006m_5nm} y \ref{fig:TvsE_11m_5nm}, correspondientes a \textbf{$m* = 0.06 m_e$} y \textbf{$m* = 1.1 m_e$}, respectivamente. \\ 
Lo primero que salta a la vista es la diferencia entre los órdenes de magnitud del eje vertical. Para un valor de 0.2eV, el primer caso muestra una transmisión de $~10^-3$, mientras que para lamisma energía, el segundo caso muestra una transmisión de $~10^-14$, con lo que hay once órdenes de magnitud entre ambos casos. Esto muestra cómo para una longitud pequeña de un transistor ($L=5$nm) la masa efectiva del material afecta de forma exponencial al comportamiento del transistor, tal y como podemos esperar según \ref{eq:tunneling}. 

Por otro lado, si observamos en el caso \ref{fig:TvsE_11m_5nm}, cuando la energía del  electrón en la fuente supera la de la barrera con el canal ($E>0.6$eV), es cuando se alcanza la transmisión máxima (como es de esperar, ya que el electrón tiene la energía suficiente para saltar la barrrera). Sin embargo, notamos que presenta oscilaciones en torno a dicho valor. Esto muestra el comportamiento de interferencia de la onda consigo misma cuando se refleja al atravesar la barrera. 


\paragraph{}Una vez identificada esta dependendia directa entre la masa efeectiva y el coeficiente de transmisión, calculamos la longitud crítica 